\originalchapter{微波背景的温度谱与宇宙学参数}
本章对应8.902第17--18讲。CMB近黑体的平均谱说明早期热史，天空温度起伏则同时记录初始扰动、光子重子流体和光传播几何。要从图中的一个峰解释参数，必须先明确温度场、角功率谱和声学尺度的定义；下述领先阶模型解释机制，不替代完整Boltzmann计算与观测似然分析。

\originalsection{收到的温度起伏包含哪些物理项}
令$T(\mathbf n)=\bar T[1+\Theta(\mathbf n)]$，$\mathbf n$指观测者看向天空的方向。最后散射附近，光子的本征温度、局部势引起的引力红移及流体本动速度都影响接收频率。用$c=1$、上一章的共形Newton规范作示意，并扣除观测者处的单极和偶极，有
\[
 \Theta_{\rm obs}\simeq[\Theta_0+\Psi-\mathbf n\cdot\mathbf v_b]_* 
 +\int_{\eta_*}^{\eta_0}(\Phi'+\Psi')\dd\eta.
\]
这是瞬时最后散射和无后续散射近似；完整式要对可见度函数和沿光线位置积分。$\Theta_0$是局部光子温度单极，$\Psi$是引力势项，速度项的符号随光线方向约定改变，积分项称积分Sachs--Wolfe效应。静态势中进入与离开时的频移相消，演化势则不能完全相消。

在物质主导、绝热长波、$\Phi=\Psi$且忽略衰减模时，Newton规范密度满足$\delta_m\simeq-2\Psi$、$\delta_\gamma=4\delta_m/3$，故$\Theta_0=\delta_\gamma/4=-2\Psi/3$，普通Sachs--Wolfe项为$\Theta_0+\Psi=\Psi/3$。势阱的本征温度与引力红移部分抵消；不能将其说成任意尺度上Sachs--Wolfe与Doppler两者自动抵消。光传播的时间延迟也是几何效应，但不能把所有时间延迟都当成上述温度抵消的充分解释。

复合后仍有多个二次效应。再电离产生新的Thomson散射，衰减原有较小尺度温度起伏的振幅，并生成大尺度偏振；散射具有角分布，不是精确各向同性。变化的势产生积分Sachs--Wolfe项，透镜改变天空方向并重新分配角功率。热Sunyaev--Zel'dovich效应是\emph{热电子对光子的逆Compton散射}，可使低频强度降低、高频升高，改变谱形；运动SZ来自电子整体速度。它们不能都当成最后散射面上的初始密度变化。

\originalsection{用球谐而不是平面波描述整个天空}
取归一球谐$\int Y_{\ell m}^*Y_{\ell'm'}\dd\Omega=\delta_{\ell\ell'}\delta_{mm'}$，有
\[
 \Theta(\mathbf n)=\sum_{\ell=0}^{\infty}\sum_{m=-\ell}^{\ell}a_{\ell m}Y_{\ell m}(\mathbf n),
 \qquad a_{\ell m}=\int\Theta Y_{\ell m}^*\dd\Omega.
\]
统计各向同性使$\langle a_{\ell m}a_{\ell'm'}^*\rangle=C_\ell\delta_{\ell\ell'}\delta_{mm'}$。对一张天空估计$\widehat C_\ell=(2\ell+1)^{-1}\sum_m|a_{\ell m}|^2$。$C_\ell$是随机场的期望谱，估计器的$m$平均不等于多个独立宇宙的平均。

$\ell=0$是平均温度，$\ell=1$的主导项来自太阳系相对CMB的运动，并有地球轨道的时变贡献；$\ell=2$为四极。特征角尺度大约$\vartheta\sim\pi/\ell$，只是数量级映射，不是精确的一一波长对应。

加法定理给角相关函数
\[
 C(\vartheta)=\langle\Theta(\mathbf n)\Theta(\mathbf n')\rangle
 =\sum_\ell\frac{2\ell+1}{4\pi}C_\ell P_\ell(\cos\vartheta).
\]
常绘$\mathcal D_\ell=\ell(\ell+1)C_\ell/(2\pi)$；若温度系数以$\mu\mathrm K$而不是$\Theta$表示，就还乘$\bar T^2$并换单位。课程图中的$\mu\mathrm K^2$不能与无量纲$C_\ell$直接相等。全空、Gaussian、无噪声时
\[
 \frac{\operatorname{Var}(\widehat C_\ell)}{C_\ell^2}=\frac{2}{2\ell+1}.
\]
这来自有限的独立$m$自由度，称宇宙方差；低$\ell$即使仪器完美也不能任意精确，掩膜和噪声则进一步改变误差。

\originalsection{为什么尺度不变谱产生大角度平台}
在普通Sachs--Wolfe近似下，曲率模投影到半径$D_M(z_*)$的天空，温度传递函数约为$\Delta_\ell(k)=j_\ell(k\chi_*)/5$，这里取平直空间和对应曲率符号约定。角谱
\[
 C_\ell=4\pi\int\mathcal P_{\mathcal R}(k)|\Delta_\ell(k)|^2\dd\ln k.
\]
若$\mathcal P_{\mathcal R}=A_s$恒定，利用$\int_0^\infty j_\ell^2(x)\dd x/x=1/[2\ell(\ell+1)]$（$\ell\ge1$），则
\[
 C_\ell\simeq\frac{2\pi A_s}{25\ell(\ell+1)},\qquad \mathcal D_\ell\simeq\frac{A_s}{25}.
\]
这解释大角度近平台。晚期积分Sachs--Wolfe、有限最后散射宽度、张量模和略倾斜初始谱会改变它，不应把低$\ell$平台当成在所有背景下精确成立。

\originalsection{声学峰与声视界的积分}
紧耦合流体的声速为$c/\sqrt{3(1+R_b)}$。忽略驱动和系数慢变，Fourier扰动满足声学振子方程，解的相位为$kr_s$，其中\emph{共动}声视界
\begin{equation}\label{eq:astro-sound-horizon}
 r_s(a_*)=\int_0^{a_*}\frac{c_s(a)\dd a}{a^2H(a)}
 =\int_{z_*}^{\infty}\frac{c_s(z)\dd z}{H(z)}.
\end{equation}
原课用$t_*c/\sqrt3$估物理声程，忽略了膨胀期间速度与尺度因子的积分；不能用于精确峰位。观察角为$\theta_s=r_s/D_M(z_*)$，基准声学多极$\ell_A=\pi/\theta_s$。各峰近乎位于整数$kr_s=n\pi$的压缩或稀疏状态，但势演化和相位移使第一温度峰\emph{不精确位于}$\ell_A$。课堂常说第一峰约$\ell\sim200$，而常用模型的$\ell_A$约300，两者指的是不同量。

绝热初态让不同波数的振子有相干初始相位，故投影后保留一系列峰。压缩与稀疏对应奇偶峰；速度项与密度项相差约四分之一周期，会改变实际温度谱。无限小尺度不能无限振荡，光子扩散产生近似$\exp[-(k/k_D)^2]$的振幅阻尼，有限最后散射宽度也平滑起伏。课程的约5角分是特定尺度估计，不是普遍锐截止。

\originalsection{几何、物质和重子各怎样影响谱}
固定早期物理组分和声视界时，空间几何主要通过$D_M$改变峰的角尺度：负曲率使相同尺看起来较小，峰向更大$\ell$移动；正曲率相反。但$H_0$、暗能量和曲率都影响距离，CMB存在几何简并，不能仅凭第一峰一项精确分别测出全部参数。

早期方程更直接依赖物理密度$\omega_m=\Omega_mh^2$、$\omega_b=\Omega_bh^2$。$\omega_m$改变辐射—物质相等时刻和势阱衰减：辐射占比更高时，变化的势驱动声学振子，产生辐射驱动和早期积分Sachs--Wolfe贡献。原课用“同一红移下年龄较小、来不及形成结构”解释峰高变化，是直观概括；其结论还依赖固定哪些参数，不能直接作为任意改变$\Omega_m$的单调律。

提高$\omega_b$增加$R_b$，改变声速和振子惯性。在近恒定势阱中，它还移动平衡位置：压缩加深，稀疏不能对称地离开原势阱，因而相对增强奇峰、抑制偶峰。第一与第二峰的比值提供重子加载信息，结合其余峰、倾斜谱和阻尼尾才可区分其他效应。BBN的重子约束提供独立交叉检验。

再电离光学深度$\tau$使原温度谱大部分乘$\ee^{-2\tau}$，与初始幅度$A_s$简并；偏振可帮助打破这一简并。透镜、前景和仪器波束也参与实际似然。本卷不把示意机制宣称为已实施完整参数估计。课程引用的Planck 2018温度谱是2020年发表的CC BY图，本卷仅保留来源说明，不复制图；这里的公式与机制为独立中文推导。
